\chapter{Design patterns}
    For each design pattern it is essential to include how we describe the same concept a design pattern explains in natural language as well
    \section{Generating Objects}
        
        \paragraph{Factory method}
            \begin{lstlisting}[language=Python]
                from abc import ABC, abstractmethod
            
                class Document(ABC):
                    @abstractmethod
                    def create_page(self):
                            pass
            
                    class Resume(Document):
                        def create_page(self):
                            return "ResumePage"
            
                    class Report(Document):
                        def create_page(self):
                            return "ReportPage"
            \end{lstlisting}
        
        \paragraph{Static factory method}
            
            \begin{lstlisting}[language=Python]
                class Car:
                    def init(self, brand):
                        self.brand = brand
                    @staticmethod
                    def toyota():
                        return Car("Toyota")
                my_car = Car.toyota()
            \end{lstlisting}
        
        \paragraph{Abstarct factory}
        
        \paragraph{Simple factory}
            
            \begin{lstlisting}[language=Python]
                def shape_factory(kind):
                    if kind == "circle":
                        return Circle()
                    elif kind == "square":
                        return Square()
            \end{lstlisting}
        
        \paragraph{Factory as a Function Object}
            
            \begin{lstlisting}[language=Python]
                class ConnectionFactory:
                    def call(self, host):
                        return Connection(host)
                factory = ConnectionFactory()
                conn = factory("localhost")
            \end{lstlisting}
        
        \paragraph{Parameterized Factory}
            
            Factory methods that accept parameters to configure the returned object.
            
            Often combined with dependency injection.
        
        \paragraph{Virtual Constructor (a.k.a. Class Method Factory)}
            
            \begin{lstlisting}[language=Python]
                class Shape:
                    @classmethod
                    def from_sides(cls, sides):
                            if sides == 3:
                                return Triangle()
                            elif sides == 4:
                                return Square()
                            \end{lstlisting}
        
        \paragraph{Factory registry}
            
            The factory registry pattern is just a way to avoid a big if/elif chain by keeping a dictionary of “keys → constructor functions”. It’s basically a lookup table for creating objects.
            
            
            
            \begin{lstlisting}[language=Python]
                class AnimalFactory:
                    _registry: dict[str, type] = {}
                    @classmethod
                    def register(cls, name: str, animal_cls: type) -> None:
                        """Register a
                    class under a given name."""
                        cls._registry[name.lower()] = animal_cls
            
                    @classmethod
                    def create(cls, name: str, *args, **kwargs):
                        """Create an instance of the registered class."""
                        animal_cls = cls._registry.get(name.lower())
                        if animal_cls is None:
                            raise ValueError(f"Unknown animal type: {name}")
                        return animal_cls(*args, **kwargs)
            
                # Example classes
                class Dog:
                    def speak(self):
                        return "woof"
            
                class Cat:
                    def speak(self):
                        return "meow"
            
                # Register them
                AnimalFactory.register("dog", Dog)
                AnimalFactory.register("cat", Cat)
                # Usage
                a = AnimalFactory.create("dog")
                b = AnimalFactory.create("cat")
                print(a.speak()) # woof print(b.speak()) # meow
            \end{lstlisting}

\section{Flexiable objects}
    
    \subsection{Compposite}
    
    \subsection{Decorator}
        \paragraph{UML}
            
        \paragraph{usecase:}
        \begin{itemize}
            \item Adding methods to a basic object
            \item Changing the behavior of a basic object
        \end{itemize}
        
        \paragraph{Implementation}
            \begin{itemize}
                    \item Interface: The genetic of what is trying to be modeled
                    \item Abstract class \textbf{Decoratr} inherits(implements) from \textbf{Interface}, this class is the genetics of being the concept of a decoration. It composes a protected \textbf{\_inner} property of type Interface. What ever method or property of it is called, it calls the same method on the \textbf{\_inner}
                    \item Decorations are the classes who inherit from \textit{Decorator}. Either they implement a method or from the Decorator or if that method doesnt exist automatically , according to inheritance rules, once again \textbf{Decorator.\_inner.method or .property} will send it to next, more inner method or property. Another option is
            \end{itemize}
        
        \paragraph{Notes}
            \begin{itemize}
                    \item Decoration reimplements inheritince in horizentally as real inheritance does it vertically. Or maybe it is better to say that it gives us a second utility for inhertitence
                    \item Decorations are not needed to implement \textbf{\_\_init\_\_} unless they neeed their own arguments for the \textbf{\_\_init\_\_(self, inner:Decorator, a, b, c)} in that case \textbf{super().\_\_init(inner)\_\_} must be called in \textbf{\_\_init\_\_} body.
                    \item In case a special kind of decorator can only decorate another special then \textbf{\_\_init\_\_(inner:SpecialDecoration, ...)}
                    \item if two decorators can not be used together then you probably need two set of decorators. or you need several sub classes from the Basic class
                    \item the docoration does not need to use \mintinline{python}{self._inner.method()}, it just calls \mintinline{python}{self.method()} and it calss
                    \item  it is always better to used passive form names for decoratars for example "diced" for something that adds dictionary abilities to inner object
            \end{itemize}
        
        \paragraph{Downsides}
            \begin{itemize}
                    \item sometime two decorators are inconsistenet with each other
            \end{itemize}
        
        \paragraph{Ralation with natural language}
            \begin{itemize}
                \item The question is if and how we describe decorated concepts in natural language
            \end{itemize}
            
        
    
    \subsection{Facade}
    
    \subsection{Template method}
    
    \subsection{Adapter}
        Makes an external SDK similar to how you see the world. Translates one interface into another so code expecting a Target can use an incompatible Adaptee without changing either.
        \begin{lstlisting}[language=Python]
            class Adaptee:
                def specific(self):
                    return "OK"

            class Adapter:
                def __init__(self, a):
                    self.a = a
                def request(self):
                    return self.a.specific()

            print(Adapter(Adaptee()).request())
        \end{lstlisting}
        A real world example
        \begin{lstlisting}[language=Python]
            from typing import Protocol

            # Target interface expected by the client
            class PaymentGateway(Protocol):
                def charge(self, amount: float, currency: str) -> str: ...

            # Adaptee (3rd-party SDK) with an incompatible API
            class StripeSDK:
                def create_payment(self, cents: int, curr: str) -> dict:
                    # Imagine real SDK call; returns a dict with an "id"
                    return {"id": f"pay_{cents}_{curr}"}

            # Adapter: makes StripeSDK look like PaymentGateway
            class StripeAdapter(PaymentGateway):
                def __init__(self, sdk: StripeSDK):
                    self._sdk = sdk

                def charge(self, amount: float, currency: str) -> str:
                    # Translate units & parameter names
                    cents = int(round(amount * 100))
                    result = self._sdk.create_payment(cents=cents, curr=currency.lower())
                    return result["id"]

            # Client code uses the Target interface only
            def checkout(gw: PaymentGateway) -> str:
                return gw.charge(12.99, "EUR")

            gateway = StripeAdapter(StripeSDK())
            payment_id = checkout(gateway)  # works via the adapter
        \end{lstlisting}

\section{Behavioral}
    \subsection{Observer}